% Bài làm buổi 20/10 — Na
\begin{problem}[8DS-PT1-004]
$\dfrac{2x - 1}{3} - \dfrac{x + 2}{4} = \dfrac{x - 1}{6}$
$\dfrac{4(2x - 1)}{12} - \dfrac{3(x + 2)}{12} = \dfrac{2(x - 1)}{12}$
$4(2x - 1) - 3(x + 2) = 2(x - 1)$
$8x - 4 - 3x - 6 = 2x - 2$
$5x - 10 = 2x - 2$
$3x = 8$
$x = \dfrac{8}{3}$
Vậy $x = \dfrac{8}{3}$.
\end{problem}

\begin{problem}[8DS-PTNT-006]
$3x^2 - 8x + 4 = 3x^2 - 6x - 2x + 4$
$= 3x(x - 2) - 2(x - 2)$
$= (x - 2)(3x - 2)$
\end{problem}

\begin{problem}[8DS-HDT-005]
\[
B = (2x + 5)^2 - 2(2x + 5)(2x - 1) + (2x - 1)^2
  = \left[(2x + 5) - (2x - 1)\right]^2
\]
$B = (2x + 5 - 2x + 1)^2 = 6^2 = 36$
Vậy $B = 36$ với mọi $x$, nên giá trị của $B$ không phụ thuộc vào $x$.
\end{problem}

\begin{problem}[8DS-HSBN-002]
a) (1) $y = 5 - 2x$ là hàm số bậc nhất với $a = -2$, $b = 5$.
(2) $y = 2x^2 + 1$ không là hàm số bậc nhất vì có $x^2$.
(3) $y = \dfrac{x}{3}$ là hàm số bậc nhất với $a = \dfrac{1}{3}$, $b = 0$.
(4) $y = 2(x - 1) - 2x + 3$ là hàm số bậc nhất với $a = 2$, $b = 1$.
b) Hàm số $y = (2m - 4)x + 1$ là hàm số bậc nhất khi $2m - 4 \neq 0$
$m \neq 2$
\end{problem}
